\chapterepigraph{Life must be understood backwards. But they forget the other
proposition, that it must be lived forwards.}{S\o ren Kierkegaard}{\textit{Journals}}

\begingroup
\let\clearpage\relax
\let\cleardoublepage\relax
\chapter{Discussion}
\label{ch:discussion}
\endgroup

\section{Geometry between architecture and optimiser}
\label{sec:repr-discussion}

The three parts of this thesis converge on one object. Architecture fixes the shape of the
variational tangent space; the optimiser acts on that space through its metric; and the
sampling measure decides how well the metric can be estimated. Read in that order, results
that arrived separately form a connected account.

Its first link is the one the energy cannot see. Two correlator families reach the same
reference-quality energy at strong confinement while carrying their Fisher weight
differently. Both the effective dimension and the leading tangent modes depend on the
parameter coordinates. Under the implemented parametrisations, the message-passing models
use fewer effective directions and are more completely described by the selected
four-operator dictionary. No random-dictionary null or held-out fit was run, so this is a
relative description, not a test of whether the remaining directions are physical. The
state-level result is simpler: the message-passing correlator has lower local-energy variance
at every confinement. That establishes that the trained states differ, although the
unmatched learned features prevent attributing the difference to message passing alone.

The second link is where the optimiser meets that space. The result there is narrower than
the usual one about natural gradients. Under $|\Psi|^2$ sampling the empirical metric is badly conditioned, $\kappa\sim10^{7}$--$10^{10}$ on precisely the models where Adam and SR agree to four decimals, so conditioning cannot be the criterion. What we propose separates the regimes
is whether a batch resolves the metric at all. Under $|\Psi|^2$ the samples come from the
measure that defines $S$. Under a fixed collocation proposal, SR is required in the tested
configurations, while CG-SR becomes unstable once proposal overlap collapses. The negative half of that is
measured. The positive half---that collocation under a fixed proposal occupies an
anisotropic but still estimable middle regime---is inferred from the estimator, since
$\kappa$ was never obtained under the collocation measure at matched batch size.

Whether the chain generalises is a fair question and the honest answer is that we cannot
settle it here. Message-passing neural quantum states have been adopted for extended fermion
systems on the argument that relational structure
scales~\cite{Pescia2024-MessagePassingNQS,KimEtAl2023-UltracoldFermiNQS}; here the tangent
diagnostics identify where the implemented families first differ. The dictionary gap widens
over the sizes measured, but three sizes do not establish a scaling law, and the comparison
is seeded at $N{=}6$.

\paragraph{What agreement with the benchmarks supports.}
\label{sec:energy-discussion}
The energies agree closely with the available fixed-node references over the benchmarked
grid, and that agreement is neither exactness nor a meeting of two blind calculations.
Three separate things limit what it shows. The first is ancestry: where a DMC value exists
it is also the target of residual pretraining. The pure-energy SR tail that follows
minimises $\langle\hat H\rangle$ under $|\Psi_\theta|^2$ and no longer sees the target, which
strengthens the comparison considerably, but the final states still descend from
reference-guided training. The second is what an energy can certify. The variational bound
keeps the reported expectation above $E_0$ whatever the training aimed at, but it says
nothing about how close to $E_0$ it lies, and a flexible ansatz can reproduce an approximate
reference without describing every other property of the exact state. Our evidence beyond
the energy is partial: the local-energy variance, which vanishes only for an eigenstate, and
a squared overlap of $1.000000$ with the exact state where one is available, at $N{=}2$. The
third concerns the cells where our energies sit below the reference. The reference is
variational only with respect to its own imposed nodes, so there is no contradiction, and
better nodes in the network could explain it. A convention difference in the Hamiltonian
would not: the exact two-electron value is reproduced, and a convention mismatch would shift
every cell. That leaves better nodes against an extrapolation or sampling error on either
side, or an error in the chain of compilations through which the reference reached us, and
we cannot tell these apart. A
fixed-node DMC calculation that uses the network's nodes would be the cleanest test. The
tables accordingly report deviation, not error, and what we claim is close agreement with the
available fixed-node references, not a validation of the wavefunctions. Below $\omega=0.1$
for $N\ge6$ we have no external value to compare with, and those states are predictions.

\section{Physics returned to the solver}
\label{sec:feedback-discussion}

A particularly useful result is that measurements from the trained state improve the
proposal used for later training.

The order matters. A generic three-component Gaussian proposal covers the strong-confinement
droplet adequately and collapses in the Wigner regime, where the effective sample size falls
to about $1.5$ of $4096$ and the gradient is dominated by a few samples. Nothing in the objective
or the architecture diagnoses that. What diagnoses it is the structure measured from the
trained states themselves: a single ring at $N{=}6$ with a well-determined radius, angular
order that appears well after radial stiffening, and shell occupancies that are read rather
than assumed. Rebuilding the proposal as a radial Gaussian at the measured shell radius times
a Dirichlet angular model lifted the effective sample size by $15$--$32\times$ with no
network training, although the absolute ESS remained low. It carried a cascade with
MCMC-free updates from $\omega\approx0.01$ to a stable checkpoint at $\omega=0.0035$.

Together these observations support a working mechanism: proposal--target mismatch produces
heavier weights, noisier gradients and less stable optimisation, while inversion of an
unreliable metric can amplify that noise. The interventions change several links at once, so
the present experiments do not establish the sequence causally. At $N{=}20$, memory limits
also force smaller models or batches and prevent the implementation from separating those
effects. A better proposal nevertheless moves the observed frontier.

Below $0.0035$ the obstacle changes character. Metropolis walkers started on the classical
shell relax inward regardless; the displacement is too small to bridge the gap; and the
collapse is non-monotonic across proposal variants. The evidence therefore points to the
warm start. At these confinements, the component decomposition shows Jastrow amplification
comparable to the determinant's suppression of the Wigner shell. After a $23\%$ change in
$\omega$, the transferred density falls back towards the non-interacting shell. This is consistent
with an imbalance between the Slater and Jastrow factors during transfer, although the
backflow unit convention remains another unresolved source of non-covariance. We have not
shown the ansatz cannot represent the $\omega=0.0027$ state, and
there is evidence that it can represent weaker traps: trained at $\omega=10^{-3}$ directly from
a pretrained checkpoint, not carried down the cascade, it comes within $0.35\%$ of our SR--VMC
energy (Table~\ref{tab:reliability}). The limit at $0.0035$ belongs to the cascade, and the
remedy we suggest, per-$\omega$ re-optimisation, possibly against a shell-amplitude-anchored
objective, is what that campaign already does.

The same failure mode may extend beyond the quantum dot: any importance-sampled scheme whose
proposal increasingly misses a strongly correlated target should monitor whether the metric
it inverts remains statistically resolved at the available sample size.

\paragraph{The staged crossover.}
The physics that drove the loop stands in a particular relation to what was known. That
parabolic dots form Wigner molecules with shell structure is established
~\cite{Egger_1999,Filinov_2001,Kong_2002,manninen2007metalclustersquantumdots,Ghosal2007-WignerMolecules}.
Nor is the separation of radial and angular ordering new: earlier simulations and
symmetry-restored calculations report radial order before orientational correlations
~\cite{Filinov_2001,Cavaliere2009-WignerTransport}. The contribution here is a
topology-conditioned measurement within the retained fixed-$S_z$ states, using finite-$n$
nulls and a held-out shell reconstruction across the selected grid. The $N{=}6$ ring follows
the classical $1.334\,\omega^{-2/3}$ scaling and reaches within $1.7\%$ at the weakest trap.
Its bond order remains near the random-angle null through $\omega=0.1$ and becomes large only
at $\omega=10^{-3}$. Whether the same staging
survives a total-spin scan we cannot say.

\paragraph{What ``MCMC-free'' means here.}
\label{sec:collocation-discussion}
The claim is about the gradient: no configuration entering a parameter update comes from a
Markov chain. A small Metropolis probe runs outside the gradient loop to select checkpoints
and certify the final energy, and no sample it produces reaches an update. It is not a claim
of freedom from Monte Carlo, since importance sampling and resampling are themselves Monte
Carlo; nor of freedom from the reference, since where a DMC value exists the residual
pretraining anneals toward it and the rollback monitors against it. Nor, at the deep-Wigner
rungs where we have no reference, is the state reference-free in ancestry: its warm starts
descend from benchmark-guided states at stronger confinement. Against the SR route,
collocation is several times less accurate across the benchmarked range---within about a
tenth of a percent of DMC at strong confinement and three tenths at $N{=}6$,
$\omega{=}0.1$---and not competitive at $N{=}20$. The two look complementary rather
than rival: collocation for fast, parallel exploration and pretraining with MCMC-free
gradient updates; SR for
final accuracy where correlation is strong.

\section{Scope, and what to do next}
\label{sec:implications}

Five limits bound everything above. All states are the lowest found within a fixed
spin-projection sector, and total spin is not scanned; that matters most in the deep-Wigner
regime, where spin transitions are plausible and the structural conclusions are strongest.
The architecture comparison is seeded at $N{=}6$, rests on fewer models at $N{=}12$ and on a
single production state at $N{=}20$, and the error bars on the production energies do not
include the variation between independent training runs. The backflow variants are not
matched in capacity, input scaling or centre-of-mass handling, so neither the
cross-architecture energy gap nor the rank collapse isolates architecture. The correlators
also use different primitive pair features. The tangent-space quantities are properties of
the networks as parametrised, not of the states they represent. And one entry in the
backflow rank table, at $N{=}20$ and $\omega=0.01$, comes from a re-run whose diagnostics
were not retained.

Five things follow that we would do next.

\emph{Make the proposal learned, not designed.} Up to the frontier the proposal is the
bottleneck and not the objective, which argues for a better sampler rather than a better
loss: normalising flows trained against $|\Psi|^2$, or the shell-informed construction of
\S\ref{sec:kernel-q3-frontier} generalised beyond ring geometries. Past the frontier the
target changes: what is needed there is an $\ell$-covariant parameterisation, or an objective
anchored on the shell amplitude, so that a stable state survives being carried to a weaker
trap.

\emph{Make rotation equivariance exact.} Permutation equivariance holds by construction and
the production CTNN projects out the centre of mass, but its displacement head reads a $d$-vector from a
network that has seen absolute coordinates, so $\mathrm{SO}(2)$ equivariance is only
encouraged---by relative-vector edge features and by the random in-plane rotation applied to
every batch---and never enforced. In a rotationally invariant trap whose weak-confinement
physics is a freely rotating molecule, that is the wrong symmetry to leave approximate.
The fix is local to one module: build the displacement as
$\Delta\mathbf r_i=\sum_{j\neq i} g_\beta(\text{invariants}_{ij})\,\mathbf r_{ij}$, a learned
scalar gain on each relative vector, so the output rotates exactly with the input and the
augmentation can be dropped. We did not run it.

\emph{Break the quadratic backflow.} At twenty electrons the tested implementation is limited
by the $\mathcal{O}(N^2 h_{\rm bf})$ edge tensor. Sparse
or low-rank message passing, a local backflow without all-to-all edge features, or the
sample-side reductions of minSR and the kernel
formulation~\cite{ChenHeyl2024-minSR,Rende2024-KernelSR} would extend the reach without
giving up the nodal correction that is the reason for having a backflow at all.

\emph{Match the backflow comparison.} A widened conventional backflow at $182$k parameters,
with the same trap-scaled inputs, displacement unit and centre-of-mass projection, is the
control needed to turn the rank-collapse observation into an architecture ablation.

\emph{Measure the metric where the hypothesis lives.} The estimability account of
\S\ref{sec:kernel-cond} rests on a spectrum measured under $|\Psi|^2$ and inferred under
collocation. Measuring the QGT under the collocation measure at matched batch size, with the
stability of its leading eigenvalues across independent batches, a bootstrap on the
resolvable part, and its sensitivity to batch size and regularisation, would turn that
account from a hypothesis into a criterion or refute it.
